\section{Analysis}\label{sec:analysis}

This section asks where the two classifiers fail and what the 8B model changes. The
error analysis uses the final DeBERTa system's 10-seed ensemble with argmax and no dev
tuning (\S\ref{sec:encoder}); the comparison uses the paired 10-seed runs of
\S\ref{sec:llm}. All figures are on dev.

\paragraph{The model, not label noise, fails.}
Of the ensemble's 308 predictions, 137 (44.5\%) fell outside the item's reference set,
and 53 of these were on items where all three annotators agreed. Annotator disagreement
in this task is real and informative \citep{uma2021learning}, but it does not account for
the gap. Scored as the test set is scored, each annotator against the other two reaches
0.643--0.766 on \Stwo{} (mean 0.684); the same ensemble, on the same three two-annotator
reference sets, reaches 0.381, 0.380 and 0.381.

\paragraph{Commitment, not only topic.}
By TeleAI's definition (an error that involves Dodging, General or Deflection in the
prediction or the reference set), 86.9\% of the ensemble's errors are ``triangle'' errors,
against 77.7\% for their system \citep{shi-etal-2026-teleai}. Only 29.2\% of the errors,
however, have both the prediction and a reference label inside the triangle. The largest
error pairs (reference label $\rightarrow$ prediction) are Implicit $\rightarrow$ Explicit
(25), General $\rightarrow$ Implicit (20) and General $\rightarrow$ Explicit (20): they
concern how committed an answer is rather than whether it stays on topic. Length also matters. The 20 inputs still cut at 1,024 tokens are
in-set only 20\% of the time (0.580 for the other 288), and the shorter half of dev scores
0.400 on \Stwo{} against 0.283 for the longer half. Test inputs (sub-question, question and answer together) are much shorter (median
137 tokens against 474 on dev), so we infer, without being able to measure it, that dev
understates test performance on this axis; the two-annotator test references cut the
other way.

A typical error shows the knowledge an answer can depend on. Asked \emph{``Would you
campaign against Senator Joe Lieberman on Iraq?''}, the speaker replied \emph{``I'm going
to stay out of Connecticut.''} The annotators gave Dodging, Dodging and Implicit; the
ensemble predicted General ($p=0.37$). Our reading, which is an inference rather than a
measurement, is that recognising a near-answer here requires knowing that Lieberman is a
senator from Connecticut. Errors of this kind motivated the knowledge hypothesis tested
with the 8B classifier.

\paragraph{Where the 8B classifier gains.}
Before the runs, we predicted that Qwen's gain would fall mostly on the commitment boundary
(Explicit $\leftrightarrow$ Implicit, General $\leftrightarrow$ Implicit/Explicit) and more
on agreed items than on contested ones. Table~\ref{tab:analysis-classes} refutes both
parts. Per-class F1 (single models, mean over the 10 paired seeds) rises in 7 of 9
classes, most in Claims ignorance (+0.403), Declining (+0.220), Dodging (+0.103) and
Deflection (+0.089). The commitment classes move least: General +0.059, Explicit +0.057,
Implicit +0.040. Partial/half-answer is never predicted by either model, and Clarification
falls by 0.144 on only 4 dev items. The in-set rate rises by +0.076 on unanimous items,
+0.105 on items with two distinct labels and +0.048 on items with three, so the gain is not
larger on agreed items. The largest gains are on the Non-Reply classes, which turn on
recognising what kind of statement an answer is (a claim of ignorance, a refusal); General
remains the weakest frequent class (F1 0.322).

\begin{table}[t]
\centering
\small
\begin{tabular}{@{}lccc@{}}
\toprule
 & DeBERTa & Qwen & $\Delta$ \\
\midrule
\multicolumn{4}{@{}l}{\emph{per-class F1}} \\
Explicit            & 0.672 & 0.729 & +0.057 \\
Implicit            & 0.403 & 0.443 & +0.040 \\
Dodging             & 0.484 & 0.587 & +0.103 \\
General             & 0.263 & 0.322 & +0.059 \\
Deflection          & 0.287 & 0.377 & +0.089 \\
Partial/half-answer & 0.000 & 0.000 & +0.000 \\
Declining to answer & 0.381 & 0.601 & +0.220 \\
Claims ignorance    & 0.296 & 0.700 & +0.403 \\
Clarification       & 0.667 & 0.524 & $-$0.144 \\
\midrule
\multicolumn{4}{@{}l}{\emph{in-set rate by distinct annotator labels}} \\
one ($n=125$)   & 0.529 & 0.605 & +0.076 \\
two ($n=150$)   & 0.493 & 0.598 & +0.105 \\
three ($n=33$)  & 0.688 & 0.736 & +0.048 \\
\bottomrule
\end{tabular}
\caption{Where the backbone gain falls: DeBERTa-v3-large (full question, 16 epochs) and
Qwen3-8B-Base with LoRA, single models, mean over the same 10 seeds, dev \Stwo{}
(multi-reference). The pre-registered prediction of a gain concentrated on the
commitment boundary and on agreed items is refuted.}
\label{tab:analysis-classes}
\end{table}

\paragraph{Ranking, calibration and the decision rule.}
Both models rank better than they decide. For Qwen (seeds 0--2), an acceptable label is
the top prediction for 60.4\% of items, in the top two for 83.1\%, in the top three for
93.1\% and in the top five for 99.1\%. Confidence tracks correctness: across fifths of dev
from least to most confident, the in-set rate rises 0.40, 0.48, 0.56, 0.70, 0.87. The
DeBERTa ensemble shows the same pattern at a lower level, with an acceptable label in its
top three for 93.8\% of items and an in-set rate from 0.435 in the least confident fifth to
0.823 in the most confident. Most remaining errors are therefore a choice among a few
candidates the model already ranks highly.

Qwen's raw predictions are also skewed towards frequent classes. On seeds 0--2 it predicted
General for only 24 items, although General is in 113 reference sets, while it predicted
Explicit 111 times against 115. This is consistent with the effect of logit adjustment
\citep{menon2021longtail}, which corrects the class prior with a single fitted parameter:
it adds +0.062 per Qwen model (0.476 to 0.538) but only +0.005 per DeBERTa model (0.384 to
0.389). Finally, the two models do not complement each other. A probability mix, with the
weight and $\tau$ fitted inside the nested cross-validation, scores 0.543 against 0.549 for
Qwen alone (both means over 10 CV splits), and the folds place a mean weight of 0.79 on
Qwen. The system-level comparison is in Table~\ref{tab:llm-main}.
